\documentclass[11pt]{article}
\usepackage{amsmath,amssymb,amsthm}
\usepackage{algorithm}
\usepackage{algpseudocode}
\usepackage{hyperref}
\usepackage{geometry}
\geometry{margin=1in}
\newtheorem{assumption}{Assumption}
\newtheorem{theorem}{Theorem}
\newtheorem{lemma}{Lemma}
\newcommand{\E}{\mathbb{E}}
\newcommand{\R}{\mathbb{R}}

\begin{document}

\section*{Algorithm Name}
\textbf{Stochastic Gradient Descent with Warm Restarts (SGDR)}

\section*{Mathematical Formulation}
Let $f:\R^{d}\rightarrow\R$ be the objective function and let $g_t$ denote a stochastic gradient
\[
g_t=\nabla f(w_t; \xi_t),\qquad \xi_t\sim\mathcal{D},
\]
where $\mathcal{D}$ is the data distribution.  
The SGDR update is
\begin{align}
w_{t+1}&=w_t-\eta_t g_t, \label{eq:update}\\[4pt]
\eta_t&=
\begin{cases}
\eta_{\max}, & t\equiv 0\pmod M,\\[4pt]
\gamma\,\eta_{t-1}, & \text{otherwise},
\end{cases}\qquad 0<\gamma<1,\; M\in\mathbb{N}. \label{eq:schedule}
\end{align}
Thus the step size is reset to the maximal value $\eta_{\max}$ every $M$ iterations and is geometrically decayed with factor $\gamma$ between resets.

\section*{Pseudocode}
\begin{algorithm}[H]
\caption{SGDR}
\begin{algorithmic}[1]
\State \textbf{Input:} initial point $w_0\in\R^{d}$, maximal step size $\eta_{\max}>0$, decay factor $\gamma\in(0,1)$, restart period $M\in\mathbb N$, total iterations $T$.
\State Set $\eta_0\gets\eta_{\max}$.
\For{$t=0,1,\dots,T-1$}
    \State Sample a minibatch $\xi_t$ and compute stochastic gradient $g_t=\nabla f(w_t;\xi_t)$.
    \State Update parameters: $w_{t+1}\gets w_t-\eta_t g_t$ \Comment{Eq.~\eqref{eq:update}}
    \If{$(t+1)\bmod M =0$}
        \State $\eta_{t+1}\gets\eta_{\max}$ \Comment{restart}
    \Else
        \State $\eta_{t+1}\gets\gamma\,\eta_{t}$ \Comment{geometric decay}
    \EndIf
\EndFor
\State \textbf{Output:} $\{w_t\}_{t=0}^{T}$ (or optionally $\bar w_T=\frac1T\sum_{t=0}^{T-1}w_t$).
\end{algorithmic}
\end{algorithm}

\section*{Dry Run Example}
Consider the simple quadratic $f(w)=(w-3)^2$ with exact gradients (no stochasticity).  
Parameters: $w_0=0$, $\eta_{\max}=0.1$, $\gamma=0.9$, $M=5$ (restart never triggers in three steps).  

\begin{tabular}{c|c|c|c}
$t$ & $g_t=\nabla f(w_t)$ & $\eta_t$ & $w_{t+1}=w_t-\eta_t g_t$ \\ \hline
0 & $2(w_0-3)=-6$ & $0.1$ & $w_1=0-0.1(-6)=0.6$ \\ 
1 & $2(w_1-3)=2(0.6-3)=-4.8$ & $0.09$ & $w_2=0.6-0.09(-4.8)=1.032$ \\ 
2 & $2(w_2-3)=2(1.032-3)=-3.936$ & $0.081$ & $w_3=1.032-0.081(-3.936)=1.350$ 
\end{tabular}

Objective values:
\[
f(w_0)=9,\; f(w_1)=(0.6-3)^2=5.76,\; f(w_2)=(1.032-3)^2\approx3.86,\; f(w_3)=(1.350-3)^2\approx2.74.
\]

The iterates move toward the minimiser $w^\star=3$, illustrating the decay of the step size.

\newpage
\section*{Assumptions}
\begin{assumption}[Smoothness]\label{ass:smooth}
$f$ is $L$‑smooth, i.e.
\[
\|\nabla f(x)-\nabla f(y)\|\le L\|x-y\|,\qquad\forall x,y\in\R^{d}.
\]
Equivalently,
\[
f(y)\le f(x)+\langle\nabla f(x),y-x\rangle+\frac{L}{2}\|y-x\|^{2}.
\]
\end{assumption}

\begin{assumption}[Unbiased Stochastic Gradient]\label{ass:unbiased}
For every $t$, conditioned on the past,
\[
\E[g_t\mid\mathcal{F}_t]=\nabla f(w_t),
\]
where $\mathcal{F}_t=\sigma(w_0,\dots,w_t,\xi_0,\dots,\xi_{t-1})$.
\end{assumption}

\begin{assumption}[Bounded Variance]\label{ass:var}
There exists $\sigma^{2}>0$ such that
\[
\E\big[\|g_t-\nabla f(w_t)\|^{2}\mid\mathcal{F}_t\big]\le\sigma^{2},\qquad\forall t.
\]
\end{assumption}

\begin{assumption}[Bounded Gradients]\label{ass:gradbound}
There is $G>0$ with $\|\nabla f(w)\|\le G$ for all $w$ in the domain explored by the algorithm.
\end{assumption}

All constants $L,G,\sigma$ are independent of $t$ and of the restart schedule.

\section*{Main Convergence Theorem}
\begin{theorem}[Non‑convex SGDR convergence]\label{thm:sgdr}
Under Assumptions~\ref{ass:smooth}–\ref{ass:gradbound} and for any total iteration budget $T$, let $\{\eta_t\}_{t=0}^{T-1}$ be generated by the warm‑restart schedule \eqref{eq:schedule}. Define
\[
\bar\eta_T\;:=\;\frac{1}{T}\sum_{t=0}^{T-1}\eta_t,\qquad
\bar\eta_T^{(2)}\;:=\;\frac{1}{T}\sum_{t=0}^{T-1}\eta_t^{2}.
\]
Then the iterates $\{w_t\}$ produced by SGDR satisfy
\begin{align}
\frac{1}{T}\sum_{t=0}^{T-1}\E\big[\|\nabla f(w_t)\|^{2}\big]
\;\le\;
\frac{2\big(f(w_0)-f^{\star}\big)}{L\,\bar\eta_T}
\;+\;
\frac{L\,\sigma^{2}\,\bar\eta_T^{(2)}}{\bar\eta_T},
\tag{1}
\end{align}
where $f^{\star}=\inf_{w}f(w)$.  

In particular, if $M$ and $\gamma$ are fixed and $T$ is a multiple of $M$, then
\[
\bar\eta_T\ge \frac{\eta_{\max}}{M}\,\frac{1-\gamma^{M}}{1-\gamma},
\qquad
\bar\eta_T^{(2)}\le \frac{\eta_{\max}^{2}}{M}\,\frac{1-\gamma^{2M}}{1-\gamma^{2}},
\]
which yields the explicit rate
\[
\frac{1}{T}\sum_{t=0}^{T-1}\E\big[\|\nabla f(w_t)\|^{2}\big]
\;=\;O\!\Big(\frac{1}{T}\Big).
\tag{2}
\]
\end{theorem}

\section*{Theoretical Properties}
\begin{itemize}
\item \textbf{Stability.} Because $\eta_t\le\eta_{\max}$ for all $t$, the standard smoothness bound guarantees that the function values cannot increase by more than $L\eta_{\max}^{2}/2$ at any iteration. The periodic restart prevents the step size from vanishing, ensuring a non‑degenerate descent direction throughout training.
\item \textbf{Hyper‑parameter Sensitivity.} The bound \eqref{eq:schedule} shows that larger $\eta_{\max}$ or longer restart period $M$ increase $\bar\eta_T$ (the denominator of (1)), thus improving the rate but potentially violating the smoothness condition if $\eta_{\max}>2/L$. The decay factor $\gamma$ controls the trade‑off between exploration (larger $\gamma$) and exploitation (smaller $\gamma$). The explicit constants in (1) make this trade‑off quantifiable.
\item \textbf{Comparison to plain SGD.} For a constant step size $\eta$, the classic non‑convex SGD bound is
\[
\frac{1}{T}\sum_{t=0}^{T-1}\E\big[\|\nabla f(w_t)\|^{2}\big]\le
\frac{2(f(w_0)-f^{\star})}{\eta T}+\frac{L\sigma^{2}\eta}{2}.
\]
SGDR replaces the constant $\eta$ by the average $\bar\eta_T$ while the second term involves $\bar\eta_T^{(2)}$. Because $\bar\eta_T^{(2)}\le\bar\eta_T\,\eta_{\max}$, the variance term is never larger than that of the best constant step size, while the first term may be smaller thanks to larger average step sizes early in each cycle.
\end{itemize}

\section*{Discussion}
The theorem demonstrates that warm restarts do not deteriorate the standard $O(1/T)$ convergence guarantee for the squared gradient norm in the non‑convex smooth setting. Moreover, the restart mechanism can improve practical performance by periodically injecting a larger learning rate, which is reflected in a larger $\bar\eta_T$ compared with a monotonically decaying schedule. The proof relies only on the classical smoothness inequality and the unbiasedness/variance assumptions; no convexity is required.

\newpage
\section*{Preliminaries (Definitions and Lemmas)}
\begin{lemma}[Smoothness inequality]\label{lem:smooth}
If $f$ is $L$‑smooth, then for any $x,y\in\R^{d}$,
\[
f(y)\le f(x)-\eta\langle\nabla f(x),g\rangle+\frac{L\eta^{2}}{2}\|g\|^{2},
\]
where $g$ is any vector and $\eta>0$.
\end{lemma}
\begin{proof}
Apply the definition of $L$‑smoothness with $y=x-\eta g$:
\[
f(x-\eta g)\le f(x)+\langle\nabla f(x),-\eta g\rangle+\frac{L}{2}\|\eta g\|^{2},
\]
which after rearranging gives the claim. \hfill $\square$
\end{proof}

\begin{lemma}[Bound on accumulated step sizes]\label{lem:averages}
For the schedule \eqref{eq:schedule} with period $M$ and decay $\gamma\in(0,1)$,
\[
\sum_{k=0}^{M-1}\eta_{t+k}
= \eta_{\max}\sum_{j=0}^{M-1}\gamma^{j}
= \eta_{\max}\,\frac{1-\gamma^{M}}{1-\gamma}.
\tag{3}
\]
Moreover,
\[
\sum_{k=0}^{M-1}\eta_{t+k}^{2}
= \eta_{\max}^{2}\sum_{j=0}^{M-1}\gamma^{2j}
= \eta_{\max}^{2}\,\frac{1-\gamma^{2M}}{1-\gamma^{2}}.
\tag{4}
\]
\end{lemma}
\begin{proof}
The schedule is geometric between two restarts, thus the sums are geometric series. \hfill $\square$
\end{proof}

\section*{Main Proof (Step‑by‑Step Derivation)}
We prove Theorem~\ref{thm:sgdr}. Throughout the proof, $\mathcal{F}_t$ denotes the $\sigma$‑algebra generated by all randomness up to iteration $t$.

\noindent\textbf{[Step 1]} Apply Lemma~\ref{lem:smooth} with $x=w_t$, $y=w_{t+1}=w_t-\eta_t g_t$, and $g=g_t$:
\[
f(w_{t+1})\le f(w_t)-\eta_t\langle\nabla f(w_t),g_t\rangle+\frac{L\eta_t^{2}}{2}\|g_t\|^{2}.
\tag{5}
\]

\noindent\textbf{[Step 2]} Take conditional expectation w.r.t.\ $\mathcal{F}_t$ and use Assumption~\ref{ass:unbiased}:
\[
\E\big[f(w_{t+1})\mid\mathcal{F}_t\big]
\le f(w_t)-\eta_t\|\nabla f(w_t)\|^{2}
+\frac{L\eta_t^{2}}{2}\E\big[\|g_t\|^{2}\mid\mathcal{F}_t\big].
\tag{6}
\]

\noindent\textbf{[Step 3]} Expand $\|g_t\|^{2}$ using variance decomposition:
\[
\|g_t\|^{2}= \|\nabla f(w_t)\|^{2}
+ \|g_t-\nabla f(w_t)\|^{2}
+2\langle\nabla f(w_t),g_t-\nabla f(w_t)\rangle.
\]
Taking conditional expectation, the cross term vanishes (unbiasedness) and Assumption~\ref{ass:var} yields
\[
\E\big[\|g_t\|^{2}\mid\mathcal{F}_t\big]
\le \|\nabla f(w_t)\|^{2}+\sigma^{2}.
\tag{7}
\]

\noindent\textbf{[Step 4]} Substitute (7) into (6):
\[
\E\big[f(w_{t+1})\mid\mathcal{F}_t\big]
\le f(w_t)-\eta_t\|\nabla f(w_t)\|^{2}
+\frac{L\eta_t^{2}}{2}\big(\|\nabla f(w_t)\|^{2}+\sigma^{2}\big).
\tag{8}
\]

\noindent\textbf{[Step 5]} Rearrange (8) to isolate $\|\nabla f(w_t)\|^{2}$:
\[
\Big(\eta_t-\frac{L\eta_t^{2}}{2}\Big)\|\nabla f(w_t)\|^{2}
\le f(w_t)-\E\big[f(w_{t+1})\mid\mathcal{F}_t\big]
+\frac{L\eta_t^{2}\sigma^{2}}{2}.
\tag{9}
\]

\noindent\textbf{[Step 6]} Since $\eta_t\le\eta_{\max}\le\frac{2}{L}$ (a standard step‑size restriction for smooth non‑convex SGD) we have
\[
\eta_t-\frac{L\eta_t^{2}}{2}\ge\frac{\eta_t}{2}.
\tag{10}
\]

\noindent\textbf{[Step 7]} Using (10) in (9) gives
\[
\frac{\eta_t}{2}\,\|\nabla f(w_t)\|^{2}
\le f(w_t)-\E\big[f(w_{t+1})\mid\mathcal{F}_t\big]
+\frac{L\eta_t^{2}\sigma^{2}}{2}.
\tag{11}
\]

\noindent\textbf{[Step 8]} Take total expectation and sum over $t=0,\dots,T-1$:
\[
\frac{1}{2}\sum_{t=0}^{T-1}\E[\eta_t\|\nabla f(w_t)\|^{2}]
\le \E[f(w_0)]-\E[f(w_T)]
+\frac{L\sigma^{2}}{2}\sum_{t=0}^{T-1}\E[\eta_t^{2}].
\tag{12}
\]

\noindent\textbf{[Step 9]} Drop the non‑negative term $\E[f(w_T)]\ge f^{\star}$ and divide by $T$:
\[
\frac{1}{T}\sum_{t=0}^{T-1}\E[\|\nabla f(w_t)\|^{2}]
\le
\frac{2\big(f(w_0)-f^{\star}\big)}{T\,\bar\eta_T}
+\frac{L\sigma^{2}\,\bar\eta_T^{(2)}}{\bar\eta_T},
\tag{13}
\]
where $\bar\eta_T=\frac{1}{T}\sum_{t=0}^{T-1}\E[\eta_t]$ and
$\bar\eta_T^{(2)}=\frac{1}{T}\sum_{t=0}^{T-1}\E[\eta_t^{2}]$.
Equation (13) is exactly the bound (1) stated in Theorem~\ref{thm:sgdr}.

\noindent\textbf{[Step 10]} Compute explicit lower/upper bounds on $\bar\eta_T$ and $\bar\eta_T^{(2)}$ using Lemma~\ref{lem:averages}.  
Assume $T=K M$ for some integer $K$ (the general case follows by a truncation argument). Over each cycle the sums (3)–(4) hold, thus
\[
\sum_{t=0}^{T-1}\eta_t
=K\;\eta_{\max}\frac{1-\gamma^{M}}{1-\gamma},
\qquad
\sum_{t=0}^{T-1}\eta_t^{2}
=K\;\eta_{\max}^{2}\frac{1-\gamma^{2M}}{1-\gamma^{2}}.
\]
Dividing by $T=KM$ yields the claimed expressions for $\bar\eta_T$ and $\bar\eta_T^{(2)}$.

\noindent\textbf{[Step 11]} Plug these bounds into (13). Since both fractions are $\Theta(1)$ (they depend only on the fixed hyper‑parameters $M,\gamma,\eta_{\max}$), the right‑hand side scales as $O(1/T)$. This gives the explicit rate (2).

All steps respect the assumptions; no convexity is invoked. \hfill $\square$

\section*{Rate Derivation (With Constants)}
Define the constant
\[
C_1:=\frac{2\big(f(w_0)-f^{\star}\big)}{L}\,
\frac{1-\gamma}{\eta_{\max}\big(1-\gamma^{M}\big)},
\qquad
C_2:=\frac{\sigma^{2}}{2}\,
\frac{1-\gamma^{2M}}{(1-\gamma^{2})}\,
\frac{1-\gamma}{1-\gamma^{M}}.
\]
Using the bounds from Lemma~\ref{lem:averages} we obtain from (13)
\[
\frac{1}{T}\sum_{t=0}^{T-1}\E\big[\|\nabla f(w_t)\|^{2}\big]
\;\le\;
\frac{C_1}{T}+C_2\;\frac{1}{T}.
\]
Thus
\[
\boxed{\displaystyle
\frac{1}{T}\sum_{t=0}^{T-1}\E\big[\|\nabla f(w_t)\|^{2}\big]
\le \frac{C_1+C_2}{T}=O\!\Big(\frac{1}{T}\Big).}
\]

\section*{Special Cases}
\begin{itemize}
\item \textbf{No restart ($M=1$).} The schedule reduces to a simple geometric decay $\eta_t=\eta_{\max}\gamma^{t}$. Lemma~\ref{lem:averages} gives $\bar\eta_T\approx \eta_{\max}(1-\gamma)/(\!1-\gamma) =\eta_{\max}$ for large $T$, and the bound collapses to the standard constant‑step SGD bound with $\eta=\eta_{\max}$.
\item \textbf{Pure cosine annealing.} If we replace the geometric decay by $\eta_t=\eta_{\max}\big(1+\cos(\pi t/M)\big)/2$, the same proof applies after noting that $\eta_t\le\eta_{\max}$ and that the averages $\bar\eta_T,\bar\eta_T^{(2)}$ can be bounded by the corresponding geometric quantities (they are of the same order). Hence the $O(1/T)$ rate still holds.
\item \textbf{Deterministic gradients.} When $\sigma=0$, the second term in (13) disappears, yielding a pure $O(1/T)$ decay of the average squared gradient norm, matching the optimal rate for smooth non‑convex deterministic optimization.
\end{itemize}

\end{document}